\section{Requirements Analysis}\label{requirements-analysis}

The web application translates the MycoAI Retrieval model into a governed product for two actor classes: \textbf{User} and \textbf{Data Owner}.

\subsection{Actor Model and Permissions}\label{actor-model-and-permissions}

\begin{table}[H]
\centering
\begin{tabular}{@{}p{3cm} p{4.5cm} p{6.5cm}@{}}
\toprule
\textbf{Actor} & \textbf{Role Value} & \textbf{Capabilities} \\
\midrule
\textbf{User} & \texttt{user} & Upload images, retrieve species predictions, review/edit bounding boxes, download batch results, submit feedback and contribution proposals \\
\textbf{Data Owner} & \texttt{owner} & All User capabilities plus index reference data, CRUD species/media metadata, browse dataset, review feedback, re-index \gls{qdrant}, assess candidate models, and manage users \\
\bottomrule
\end{tabular}
\caption{Actor roles and capabilities.}
\end{table}

The Data Owner role inherits all User permissions. Backend RBAC is enforced through FastAPI dependencies, while the frontend conditionally renders owner-only pages and actions.

\subsection{Use Case Diagram}\label{use-case-diagram}

\begin{figure}[H]
\centering
\includegraphics[width=\textwidth,keepaspectratio]{02_chapter3_system/ch03_usecase_diagram.png}
\caption{Use case diagram for MycoAI Retrieval.}
\label{fig:ch03-usecase-diagram}
\end{figure}

\subsection{Functional Requirements Summary}\label{functional-requirements-summary}

\begin{table}[H]
\centering
\begin{tabular}{@{}p{3cm} p{4.5cm} p{6.5cm}@{}}
\toprule
\textbf{FR ID} & \textbf{Requirement} & \textbf{Implementation Status} \\
\midrule
FR-001 & JWT authentication for protected workflows & Implemented through authentication API and role-aware frontend shell \\
FR-006 & Species retrieval from single uploaded strain images & Implemented as upload, segmentation review, vector retrieval, and ranked result display \\
FR-007 & Species retrieval from batch uploaded folders & Implemented as asynchronous batch workflow with progress/results state \\
FR-009 & Auto-segmentation with editable bounding boxes & Implemented through deterministic KMeans segmentation and frontend review; YOLO remains research alternative \\
FR-010/011 & Same-media / all-media KNN strategy & Implemented through metadata-aware Qdrant filters and fallback search scope \\
FR-013/014 & Feedback submission and contribution proposals & Implemented as post-retrieval feedback workflow \\
FR-015 & Data Owner feedback review & Implemented as owner-only review inbox \\
FR-017 & Data Owner index new reference data & Implemented as owner-only reference-data indexing workflow \\
FR-019 & CRUD Species and Media metadata & Implemented as owner-only metadata management \\
FR-021 & Browse/search/filter/group dataset & Implemented as owner-only dataset governance page \\
FR-023 & Archive/restore only, no permanent delete & Implemented through reversible archive state \\
FR-025 & In-system Qdrant re-index trigger & Implemented through model/index maintenance workflow \\
FR-028 & Candidate Model upload, evaluate, promote/reject & Implemented as owner-controlled candidate model lifecycle \\
FR-034 & Role-based API enforcement & Implemented through backend authorization and frontend route gating \\
\bottomrule
\end{tabular}
\caption{Functional requirements and implementation status.}
\end{table}

\subsection{Key Workflows}\label{key-workflows}

\begin{table}[H]
\centering
\begin{tabular}{@{}p{4cm} p{6cm} p{3.5cm}@{}}
\toprule
\textbf{Actor Goal} & \textbf{System Behavior} & \textbf{Actor} \\
\midrule
Retrieve species from image & Upload, segment, extract features, KNN search in Qdrant, display ranked results & User, Data Owner \\
Submit feedback on results & Post-retrieval correction or contribution proposal submission & User, Data Owner \\
Index new reference data & Upload, segment, extract features, write vectors to Qdrant & Data Owner \\
Manage species and media metadata & CRUD operations on species, media, strain records & Data Owner \\
Browse and govern dataset & Filter, search, group, archive/restore reference records & Data Owner \\
Review submitted feedback & Inspect and act on user feedback submissions & Data Owner \\
Maintain model and index & Trigger Qdrant re-index, upload and assess candidate models & Data Owner \\
Manage users & Account creation, role assignment, audit log inspection & Data Owner \\
\bottomrule
\end{tabular}
\caption{Key workflows mapping actor goals to system behaviors.}
\end{table}

Quality targets include retrieval response under 5~seconds after segmentation, upload validation (JPEG/PNG/TIFF, 50~MB max, $\geq$256$\times$256 pixels), JWT access token expiry of 1~hour with 30-day refresh, and Data Owner mutation audit logging. All API errors follow RFC~7807 Problem Details with field-level detail for form-level mapping.

\section{System Architecture}\label{system-architecture}

\subsection{High-Level Architecture}\label{high-level-architecture}

The system follows a decoupled client-server architecture with six containerized services coordinated by Docker Compose.

\begin{figure}[H]
\centering
\includegraphics[width=\textwidth,keepaspectratio]{02_chapter3_system/ch03_architecture.png}
\caption{High-level system architecture: React SPA frontend, FastAPI backend, PostgreSQL, Qdrant, Redis, and Celery worker.}
\label{fig:ch03-architecture}
\end{figure}

\subsection{Technology Stack}\label{technology-stack}

\begin{table}[H]
\centering
\caption{Technology stack summary}
\begin{tabularx}{\textwidth}{@{}p{3cm}p{4.2cm}X@{}}
\toprule
\textbf{Layer} & \textbf{Technology} & \textbf{Justification} \\
\midrule
\textbf{Frontend} & React 19 + Vite + TypeScript~\cite{react,vite} & SPA with fast development feedback and type safety \\
\textbf{UI Kit} & shadcn/ui + Tailwind CSS v4 & Accessible components and utility styling \\
\textbf{Data Fetching} & TanStack Query & Caching and background refetch \\
\textbf{Backend} & FastAPI (Python 3.13)~\cite{fastapi,python2024} & Async-native API framework \\
\textbf{ORM} & SQLAlchemy 2.0 async & Relational data access \\
\textbf{Vector DB} & Qdrant~\cite{qdrant} & Cosine similarity search with metadata filters \\
\textbf{Task Queue} & Celery + Redis & Background jobs and retries \\
\textbf{CV Libraries} & OpenCV, NumPy, scikit-learn~\cite{bradski2000opencv,harris2020array,pedregosa2011scikit} & Segmentation, preprocessing, and research segmentation alternatives \\
\textbf{Deployment} & Docker Compose~\cite{merkel2014docker} & Reproducible single-VM deployment \\
\bottomrule
\end{tabularx}
\end{table}

\subsection{Authentication Flow}\label{authentication-flow}

\begin{figure}[H]
\centering
\includegraphics[width=\textwidth,keepaspectratio]{02_chapter3_system/ch03_auth_sequence.png}
\caption{Authentication sequence for login, token refresh, logout, and Data Owner provisioning.}
\label{fig:ch03-auth-sequence}
\end{figure}

\subsection{Error Response Format}\label{error-response-format}

All API errors follow RFC~7807 Problem Details. Validation payloads include field-level details so the frontend can map errors back to forms.

\section{Database Design}\label{database-design}

\subsection{Entity-Relationship Diagram}\label{entity-relationship-diagram}

\begin{figure}[H]
\centering
\includegraphics[width=\textwidth,keepaspectratio]{02_chapter3_system/ch03_erd.png}
\caption{Entity-Relationship Diagram for the MycoAI Retrieval PostgreSQL schema.}
\label{fig:ch03-erd}
\end{figure}

\subsection{Key Design Decisions}\label{key-design-decisions}

\textbf{UUID primary keys:} all entity tables use UUIDs to avoid sequential enumeration and simplify distributed-safe creation.

\textbf{Soft delete pattern:} species, \gls{media}, strain, image, and segment records use archive flags instead of permanent deletion.

\textbf{JSONB for flexible state:} retrieval jobs, training progress, and audit payloads use JSONB fields where rigid relational columns would be unnecessarily restrictive.

\subsection{Data Update Status Lifecycle}\label{data-update-status-lifecycle}

\begin{figure}[H]
\centering
\includegraphics[width=\textwidth,keepaspectratio]{02_chapter3_system/ch03_data_state_machine.png}
\caption{State transition diagram for the data update status field, which drives Qdrant re-indexing decisions.}
\label{fig:ch03-data-state-machine}
\end{figure}

\subsection{Dual-Database Strategy}\label{dual-database-strategy}

The system adopts a dual-database design to separate structured governance data from high-dimensional vector search:

\begin{center}
\small
\begin{tabular}{@{}p{2.8cm} p{4cm} p{5.5cm}@{}}
\toprule
\textbf{Concern} & \textbf{PostgreSQL} & \textbf{Qdrant} \\
\midrule
Data type & structured relational records & embedding vectors \\
Query model & SQL joins and filters & cosine similarity K-NN search \\
Consistency & ACID transactions & eventually consistent upserts \\
Governance & users, audit, metadata & vector freshness tracking \\
Role & source of truth & search acceleration \\
\bottomrule
\end{tabular}
\captionof{table}{PostgreSQL and Qdrant responsibilities.}
\end{center}

\normalsize

A relational database was chosen as the primary governance store not only for transactional correctness, but also because SQL-based exploration lowers the cost of development even for AI-assisted workflows. Tabular schemas are easy to inspect, modify, and extend incrementally---new metadata columns, audit tables, and index-state flags can be added through standard migrations without rebuilding the search index. This flexibility accelerates iterative feature development and simplifies debugging when agents or humans need to inspect system state directly. Qdrant is reserved exclusively for the task where a specialised index is necessary: cosine-similarity nearest-neighbour search over embedding vectors, where relational engines cannot compete at scale.

\section{Frontend Architecture}\label{frontend-architecture}

\subsection{User Interface Scope}\label{component-architecture}

The thesis describes the interface at workflow level. The implemented application contains the major pages: Retrieve, My Feedback, Dashboard, Index New Data, Dataset, Metadata, Feedback Inbox, Model Index, User Management, and Audit Log.

\subsection{Key Page: Retrieve}\label{key-page-retrieve-species-retrieval}

The Retrieve page implements a three-step workflow:

\begin{itemize}
\item \textbf{Upload}: user selects one image or a batch and provides \gls{strain}/\gls{media} metadata.
\item \textbf{Segmentation review}: user inspects and edits predicted bounding boxes.
\item \textbf{Results}: system shows ranked species predictions with supporting nearest neighbors.
\end{itemize}

\subsection{Use-Case Flow Demonstration with Agent-Browser}\label{use-case-flow-demonstration}

To document implemented use-case behavior, screenshots were captured with \texttt{agent-browser}, a headless browser automation CLI built in Rust on Chrome DevTools Protocol. This provided repeatable visual evidence for important interaction steps across authentication, retrieval, and owner-facing governance pages.

\begin{figure}[H]
\centering
\includegraphics[width=\textwidth,keepaspectratio]{11_e2e_screenshots/e2e_login.png}
\caption{Authentication flow entry point: login page used in the end-to-end browser-driven walkthrough.}
\label{fig:ch03-e2e-login}
\end{figure}

\begin{figure}[H]
\centering
\includegraphics[width=\textwidth,keepaspectratio]{11_e2e_screenshots/e2e_dashboard.png}
\caption{Post-login dashboard overview showing dataset scale, summary cards, and operator-facing status information.}
\label{fig:ch03-e2e-dashboard}
\end{figure}

\begin{figure}[H]
\centering
\includegraphics[width=\textwidth,keepaspectratio]{11_e2e_screenshots/e2e_retrieve_upload.png}
\caption{Retrieve Species upload step with batch ZIP support and retrieval configuration controls.}
\label{fig:ch03-e2e-upload}
\end{figure}

\begin{figure}[H]
\centering
\includegraphics[width=\textwidth,keepaspectratio]{11_e2e_screenshots/e2e_segmentation.png}
\caption{Segmentation review step with editable bounding boxes before final retrieval.}
\label{fig:ch03-e2e-segmentation}
\end{figure}

\begin{figure}[H]
\centering
\includegraphics[width=\textwidth,keepaspectratio]{11_e2e_screenshots/e2e_retrieval_results.png}
\caption{Retrieval result step showing ranked species predictions, confidence status, nearest-neighbour evidence, and distribution summary.}
\label{fig:ch03-e2e-results}
\end{figure}

The browser-driven walkthrough aligns with UC-AUTH-01 and UC-RETRIEVE-01 in Appendix~B, and demonstrates that implemented page flow matches the specified interaction sequence from authentication to retrieval result review.

\subsection{Key Page: Dataset Browser}\label{key-page-dataset-browser-data-owner}

The Dataset Browser supports owner-only filtering, grouped inspection, metadata review, and reversible archive/restore actions for governed reference data.

\subsection{Key Page: Dashboard}\label{key-page-dashboard-data-owner}

The Dashboard consolidates dataset scale, species/media distribution, and system status so operators can quickly assess repository health and retrieval readiness.

\subsection{Key Page: User Management and Audit Log}\label{key-page-user-management-audit-log-data-owner}

User Management and Audit Log pages provide governance controls for account administration and traceable operational changes.

\section{Core Backend Pipelines}\label{core-backend-pipelines}

\subsection{Segmentation Pipeline}\label{segmentation-pipeline}

The deployed backend uses a deterministic KMeans-driven pipeline for preprocessing and bounding-box extraction. The research chapter evaluates YOLO segmentation as a stronger alternative for flare-heavy edge cases.

\subsection{Retrieval Pipeline}\label{retrieval-pipeline}

The retrieval orchestrator performs feature extraction, Qdrant nearest-neighbor search, aggregation across segments, and final ranking. Same-media retrieval is used when the query medium is managed; otherwise the system falls back to all-media search.

\subsection{Retrieval Query Sequence}\label{retrieval-query-sequence}

\begin{figure}[H]
\centering
\includegraphics[width=\textwidth,keepaspectratio]{02_chapter3_system/ch03_srs_retrieve_sequence.png}
\caption{High-level retrieval sequence.}
\label{fig:ch03-retrieve-sequence}
\end{figure}

\subsection{Qdrant Vector Database Design}\label{qdrant-vector-database-design}

The Qdrant collection stores per-segment points with multiple named vectors, allowing different feature extractors to be queried or compared without duplicating payload metadata.

\subsection{Indexing and Feedback Sequences}\label{indexing-feedback-sequences}

Data Owner indexing reuses upload and segmentation behavior but writes reference vectors instead of returning a species prediction. Feedback remains an extension of retrieval: Users submit corrections or contribution proposals after seeing results, and Data Owners review them before any reference data change.

\begin{figure}[H]
\centering
\includegraphics[width=\textwidth,keepaspectratio]{02_chapter3_system/ch03_srs_index_sequence.png}
\caption{Data Owner reference-data indexing sequence.}
\label{fig:ch03-index-sequence}
\end{figure}

\begin{figure}[H]
\centering
\includegraphics[width=\textwidth,keepaspectratio]{02_chapter3_system/ch03_srs_feedback_sequence.png}
\caption{Feedback submission and Data Owner review sequence.}
\label{fig:ch03-feedback-sequence}
\end{figure}

\subsection{Model and Index Maintenance}\label{model-and-index-maintenance}

The Model Index page manages re-indexing and candidate-model assessment. Heavy retraining remains an offline workflow tied to the research pipeline.

\begin{figure}[H]
\centering
\includegraphics[width=\textwidth,keepaspectratio]{02_chapter3_system/ch03_candidate_model.png}
\caption{Candidate Model promotion flow (UC-MODEL-01 Maintain Model and Index): Data Owner uploads, evaluates side-by-side against the current active model on a fixed evaluation set, and manually promotes or rejects.}
\label{fig:ch03-candidate-model}
\end{figure}

\begin{figure}[H]
\centering
\includegraphics[width=\textwidth,keepaspectratio]{02_chapter3_system/ch03_reindex_warning.png}
\caption{Re-index trigger and external retraining guidance for UC-MODEL-01: the Data Owner inspects index status including changed-data and accumulated-feedback counts, re-indexes Qdrant for stale items, and receives external retraining guidance when the change threshold is crossed.}
\label{fig:ch03-reindex-warning}
\end{figure}

\section{Deployment Architecture}\label{deployment-architecture}

\subsection{Docker Compose Topology}\label{docker-compose-topology}

\begin{figure}[H]
\centering
\includegraphics[width=\textwidth,keepaspectratio]{02_chapter3_system/ch03_deployment.png}
\caption{Production Docker Compose topology: services, network boundaries, and persistent volumes.}
\label{fig:ch03-deployment}
\end{figure}

\subsection{CI/CD Pipeline}\label{cicd-pipeline}

GitHub Actions run on every push and pull request:

\begin{itemize}
\item \textbf{Backend}: ruff check, format check, mypy (strict), then pytest.
\item \textbf{Frontend}: ESLint, TypeScript typecheck, then Vite build.
\item \textbf{Docker}: image build dry-run verification.
\end{itemize}

Deployment to production is manual:

\begin{itemize}
\item Pull updated container images.
\item Restart services with the new images.
\item Triggered only after pull request approval and merge to the production branch.
\end{itemize}

\section{Summary}\label{summary}

The implemented product combines governed metadata management, deterministic backend segmentation, vector retrieval through \gls{qdrant}, and role-based workflows for laboratory usage. Together with the research pipeline in Chapter 2, it forms a full path from experimental retrieval model to operator-facing application.
